% ============================================================
\chapter{Guía de Uso y Reproducción}
\label{ap:guia_uso}
% ============================================================

Este apéndice complementa la Sección~\ref{sec:implementacion} con los pasos
concretos para instalar \textit{kubescan}, interpretar el informe que produce
sobre un caso real, y reproducir el pipeline de entrenamiento de investigación
descrito en el Capítulo~\ref{cap:diseno}.

% -----------------------------------------------------------
\section{Instalación}
% -----------------------------------------------------------

El paquete se instala en modo editable desde la raíz del repositorio:

{\footnotesize\begin{verbatim}
pip install -e kubescan/
\end{verbatim}}

Las versiones de Python soportadas y la resolución de dependencias se verifican
en la Sección~\ref{sec:evaluacion_funcional} (RNF-3).

% -----------------------------------------------------------
\section{Uso de la interfaz de línea de comandos}
% -----------------------------------------------------------

\subsection{Análisis de un directorio de manifiestos}

El comando \texttt{kubescan scan} analiza un directorio de manifiestos \ac{YAML}
de forma recursiva. A modo de ejemplo, se analiza un directorio con dos
manifiestos: un \textit{pod} con montaje del sistema de ficheros del anfitrión
y contenedor privilegiado (\texttt{pod-privilegiado.yaml}, adaptado de
\textit{BishopFox BadPods}), y un \textit{Deployment} sin
\texttt{securityContext} ni \texttt{NetworkPolicy} definidos
(\texttt{deployment-web.yaml}, adaptado de \textit{Kubernetes Goat}):

{\footnotesize\begin{verbatim}
$ kubescan scan /tmp/manifests-ejemplo --cluster-name produccion-web
Scanning /tmp/manifests-ejemplo (5 GNN fold models loaded)...
  2 manifest(s) with workload resources found

==================================================================
  KUBESCAN - Attack-Chain Risk Report
  Cluster : produccion-web
  Path    : /tmp/manifests-ejemplo
==================================================================

  VERDICT - ATTACK_CHAIN          x  HIGH RISK - review immediately

  Ensemble score    : 0.9037
  Chain probability : 0.7104   (5-fold GNN ensemble)
  Clean probability : 0.0130
  Mean RF risk      : 0.9961
  Escape fraction   : 0.5000   (1/2 manifests have escape flags)
  Lateral fraction  : 2/2 manifests have lateral flags

  Weights: w_rf=0.337  w_gnn=0.328  w_escape=0.335

==================================================================
\end{verbatim}}

La salida reporta las tres señales de entrada al ensemble (\textit{mean RF
risk}, \textit{chain probability} y \textit{escape fraction}, ponderadas por
\(w_{\text{rf}}\), \(w_{\text{gnn}}\) y \(w_{\text{esc}}\) respectivamente,
optimizados por el \ac{GA} — ver Sección~\ref{sec:capa3_diseno}), y clasifica el
clúster mediante los umbrales fijos definidos en esa misma sección. En este
caso, la presencia de un nodo con flags de escape activas (\texttt{pod-privilegiado.yaml})
eleva el \textit{ensemble\_score} muy por encima del umbral de 0,60.

El indicador \texttt{--show-nodes} añade un desglose por manifiesto (\textit{risk
score}, tipo de nodo y flags activas); \texttt{--format json} produce el mismo
informe en formato \ac{JSON} estructurado para integración con otras
herramientas. A continuación se muestra una versión abreviada (se omiten
campos y valores para mayor claridad):

{\footnotesize\begin{verbatim}
$ kubescan scan /tmp/manifests-ejemplo --cluster-name produccion-web --format json
{
  "cluster": "produccion-web",
  "verdict": "ATTACK_CHAIN",
  "ensemble_score": 0.903702,
  "chain_probability": 0.710363,
  "mean_rf_risk": 0.99605,
  "n_manifests": 2,
  "n_escape_capable": 1,
  "manifests": [
    {
      "file": "pod-privilegiado.yaml",
      "risk_score": 0.999914,
      "escape_capable": true,
      "flags": ["TRUE_HOST_PID", "TRUE_HOST_IPC", "..."]
    },
    "..."
  ]
}
\end{verbatim}}

Los pesos del modelo se resuelven mediante la cadena descrita en la
Sección~\ref{sec:implementacion} (argumento \texttt{--checkpoints-dir}
\(\to\) variable de entorno \texttt{KUBESCAN\_CHECKPOINTS} \(\to\) enlace
simbólico); no es necesario indicarlos explícitamente si el paquete se instaló
junto con el repositorio completo.

\subsection{Análisis de un clúster en ejecución}

El comando \texttt{kubescan live} aplica el mismo \textit{pipeline} sobre el
estado actual de un clúster accesible mediante \textit{kubectl}, sin necesidad
de exportar manifiestos manualmente:

{\footnotesize\begin{verbatim}
kubescan live -n produccion
kubescan live --all-namespaces --format json
\end{verbatim}}

Este comando requiere \textit{kubectl} configurado con permisos de lectura
sobre \textit{workloads} y recursos \ac{RBAC} en el clúster objetivo.

% -----------------------------------------------------------
\section{Reproducción del pipeline de investigación}
% -----------------------------------------------------------

{\raggedright
El pipeline de entrenamiento se ejecuta mediante los objetivos definidos en el
\texttt{Makefile} de la raíz del repositorio. Los pasos de adquisición y
extracción (\texttt{research/scripts/01\_acquire/}, con
\texttt{download\_github\_manifests.py} e \texttt{ingest\_attack\_repos.py}; y
\texttt{research/scripts/02\_extract/}, con \texttt{extract\_yaml\_features.py}
y \texttt{build\_graphs.py}) son intensivos en red y no están automatizados en
el \texttt{Makefile}; se ejecutan manualmente una única vez y producen los
grafos de clúster en \texttt{research/data/graphs/*.npz}.\par}

A partir de ahí, el resto del pipeline es reproducible con dos objetivos:

{\footnotesize\begin{verbatim}
make data       # aumentacion -> cache de grafos -> splits group-aware
make reproduce  # data -> RF -> GAT (5-fold) -> ensemble GA -> eval en test
\end{verbatim}}

{\raggedright
El objetivo \texttt{reproduce} encadena, en orden, \texttt{train\_rf.py},
\texttt{train\_gnn.py} (5-fold, 300 épocas, \texttt{hidden=64},
\texttt{heads=4}, \texttt{layers=3}), \texttt{run\_ga\_ensemble.py} y
\texttt{evaluate\_test\_set.py}, todos con semilla fija (\texttt{--seed 42} por
defecto), y finaliza generando
\texttt{research/models/checkpoints/run\_manifest.json} con la procedencia de
la ejecución. El objetivo \texttt{reproduce-ablations} reentrena las variantes
de arquitectura reportadas en la Tabla~\ref{tab:arch_ablation} (\ac{GAT} de
2~capas y \ac{GCN} de 3~capas) bajo las mismas semillas e hiperparámetros de
entrenamiento.\par}
